\begin{abstract}
\noindent
Completion is ordinarily treated as the successful terminal condition of production, and incompletion as a privation measured by distance from the object it has failed to become. This essay develops a contrary account built on one definition: completion is an authority's withdrawal of recognition from the continuations of an artifact within a declared scope. Completion is then not a property of a state but a relation among a state, an authority, a time, and a scope, and it is frequently a transfer of unresolved costs from an authorized producer to less powerful successors. The argument distinguishes \emph{inert} incompletion, in which work has stopped without an intelligible means of resumption, from \emph{latent} incompletion, in which the present state remains legible, accessible, and continuable along more than one admissible path, and it distinguishes three ways in which continuations are withdrawn: epistemically, operationally, and jurisdictionally.

A core model represents an artifact by its state space, recognized continuations, interpretive information, costs, allocation of authority, and identity relation. Descriptive continuational breadth is kept separate from signed, control-indexed option value, so that the value of a retained option depends on who exercises it and for whom. The principal results show that completion can raise present fitness while contracting breadth; that an option can be valuable to its controller and harmful to those it affects; that optimal closure over a dependency structure, condensed where commitments are mutually dependent, is a maximum-weight ideal that in general leaves purpose-like components open; that suggestion systems with sufficient refusal cost make arbitrary predictions self-confirming; that revision capacity can grow while the alignment of authority with affectedness declines; and that symmetric reopening rules erode settlements under repeated challenge. Standard results from the theory of irreversible investment and the value of information are identified as such. Domain models for development, interfaces, machine assistance, administration, political authority, and ruins instantiate the core model. The normative conclusion is a principle of \emph{asymmetric completion}: close the transformations whose openness would destroy the conditions of participation, and preserve those through which future participants may revise purposes that no present authority can legitimately settle in advance.

\medskip\noindent\textbf{Keywords:} completion; option value; irreversibility; continuation; modularity; choice architecture; classification; institutional design.
\end{abstract}
